\section{Towards an NPE-based Shear Estimator}
\label{sec:shear}

In Sec.~\ref{sec:validation}, we showed that our NPE produces well-calibrated posteriors. However, what we ultimately need is an estimate of the shear that is unbiased over an ensemble of galaxies, i.e. one with small $m$ and $\mathbf{c}$ in Eq.~(\ref{eqn:linear-bias}). In this section, we explore how the NPE posteriors can be turned into such an estimate. The naive choice would be to use the posterior mean of every galaxy as the point estimate for its shear. However, that estimator is biased, since each posterior mean is pulled toward the center of the prior (see Appendix~\ref{app:posterior-means}). Instead, we can make use of the full posteriors provided by our model to give us an ensemble shear estimate with minimal bias. We show in (Sec.~\ref{ssec:stacking}) that stacking relatively few likelihoods of galaxies of the same shear can reduce the multiplicative bias to $\sim1\%$. We then verify this approach in a varying cluster shear field (Sec.~\ref{ssec:cluster}), and use the likelihoods of the individual galaxies to detect substructure in the cluster (Sec.~\ref{ssec:substructure}). Through these tests, we show that our NPE upholds the main advantages of KL, even if individual shear posteriors are biased.

\subsection{Likelihood Stacking}
\label{ssec:stacking}

Rather than reducing the posterior of each galaxy to a point estimate, we can use the entire posterior that our NPE outputs, which contains more information than the point estimate itself. If a set of $N$ galaxies share the same shear $\mathbf{g}$, the posterior of that shear given the data of all galaxies is
\begin{equation}
\label{eqn:stack}
    p(\mathbf{g} \mid \mathbf{c}_1, \dots, \mathbf{c}_N) \propto \pi(\mathbf{g}) \prod_{i=1}^{N} \mathcal{L}_i(\mathbf{g}),
\end{equation}
where $\mathcal{L}_i(\mathbf{g}) = p(\mathbf{c}_i \mid \mathbf{g})$ is the likelihood of the shear for galaxy $i$, marginalized over all the other parameters. The prior on the shear enters only once, instead of once per galaxy as when averaging posterior means. As $N$ grows, the product of likelihoods becomes much narrower than the prior, so the pull of the prior, and with it the shrinkage, asymptotes to having no effect on the joint posterior. For example, in a Gaussian likelihood model, a single likelihood of width $\sigma_L$ and a prior of width $\sigma_\pi$ give a posterior mean that responds to a fraction $\sigma_\pi^2/(\sigma_\pi^2 + \sigma_L^2)$ of the true shear. However, the product of $N$ likelihoods of width $\sigma_L$ has width $\sigma_L/\sqrt{N}$, so the response of the stacked posterior mean is $\sigma_\pi^2/(\sigma_\pi^2 + \sigma_L^2/N)$, which approaches 1 as $N$ grows. This is the Bayesian shear formalism of Bernstein \& Armstrong~\citep{BA14}, which we can apply directly to our NPE. We also explain this distinction more concretely in Appendix~\ref{app:shrinkage}.

Our NPE gives us the likelihood of each galaxy directly. Its marginal posterior of the shear, $q(\mathbf{g} \mid \mathbf{c}_i)$, approximates the posterior under the training prior, so
\begin{equation}
\label{eqn:npe-likelihood}
    \mathcal{L}_i(\mathbf{g}) \propto \frac{q(\mathbf{g} \mid \mathbf{c}_i)}{\pi_\mathrm{train}(\mathbf{g})},
\end{equation}
which is simply proportional to $q(\mathbf{g} \mid \mathbf{c}_i)$ inside the prior box, since our training prior on the shear is uniform. We estimate $q(\mathbf{g} \mid \mathbf{c}_i)$ from the posterior samples of each galaxy with a kernel density estimate (KDE). We use the KDE module as implemented in \textsc{SciPy}~\citep{SciPy}, specifying it with the full covariance of the posterior samples rather than the variance for each component, as we often have posteriors that are narrow along diagonal directions in the shear plane. Since the posterior is confined to the prior box, we reflect the samples across the edges of the box to avoid underestimating the density near the edges. We also add a small constant floor to each density, so that a single galaxy whose samples happen to miss the true shear cannot set the product to zero there. We then evaluate the product of Eq.~(\ref{eqn:stack}) on a fine grid of $\mathbf{g}$ and take its mean as our estimate of the shear.

Testing Eq.~(\ref{eqn:stack}) requires groups of galaxies that share the same shear, which the Xu3 test set does not have. We therefore use a separate benchmark of 25 shear cells, with $g_1, g_2 \in \{-0.08, -0.04, 0, 0.04, 0.08\}$. Each cell contains the same 32 galaxies, whose other parameters are drawn from the training prior, with one noise realization per galaxy, as for real observations. Since these galaxies were drawn from the training prior, we use the posterior samples without the TF reweighting. In each cell, we stack the first $N$ galaxies, and measure $m$ and $c$ with a linear fit of the 25 cell estimates against the true shears of the cells. We estimate the uncertainties by bootstrap resampling the galaxies, using the same resampled galaxies in all cells. For a controlled comparison, we also average the posterior means of the same galaxies in each cell.

\begin{figure}
    \includegraphics[width=\linewidth]{figs/stacking.pdf}
    \caption{Multiplicative bias against the number of galaxies $N$ per shear cell, for 25 cells of constant shear with one noise realization per galaxy. Blue: average of the posterior means of the individual galaxies. Purple: mean of the stacked posterior of Eq.~(\ref{eqn:stack}), with the likelihood of each galaxy estimated from its NPE posterior using Eq.~(\ref{eqn:npe-likelihood}). Filled circles are $m_1$, open squares $m_2$. Markers are the medians and error bars the 16th--84th percentiles of 1000 bootstrap resamples of the galaxies.}
    \label{fig:stacking}
\end{figure}

\begin{table}
\centering
\caption{Multiplicative and additive bias of the stacked posterior of Eq.~(\ref{eqn:stack}) against the number of galaxies $N$ per shear cell, for the benchmark of Fig.~\ref{fig:stacking}. Values are the medians and errors the 16th--84th percentiles of 1000 bootstrap resamples of the galaxies.}
\label{tab:stacking}
\small
\setlength{\tabcolsep}{3.5pt}
\begin{tabular}{ccccc}
\toprule
$N$ & $10^2 m_1$ & $10^2 m_2$ & $10^4 c_1$ & $10^4 c_2$ \\
\midrule
1 & $-40.9^{+38.0}_{-44.5}$ & $-34.6^{+33.7}_{-42.6}$ & $1.5^{+98.6}_{-53.8}$ & $-2.3^{+89.6}_{-128.7}$ \\
2 & $-6.0^{+5.4}_{-33.5}$ & $-3.0^{+3.6}_{-27.3}$ & $2.4^{+64.9}_{-32.8}$ & $0.5^{+35.8}_{-39.6}$ \\
4 & $-1.3^{+1.4}_{-4.1}$ & $-0.1^{+1.4}_{-3.9}$ & $2.6^{+17.5}_{-15.0}$ & $1.4^{+13.7}_{-12.0}$ \\
8 & $-1.1^{+0.9}_{-1.1}$ & $0.1^{+0.9}_{-0.9}$ & $2.9^{+9.3}_{-11.4}$ & $2.3^{+11.3}_{-5.9}$ \\
16 & $-1.0^{+0.6}_{-0.7}$ & $0.1^{+0.6}_{-0.5}$ & $2.4^{+5.7}_{-7.0}$ & $3.6^{+7.8}_{-3.8}$ \\
32 & $-1.0^{+0.4}_{-0.5}$ & $0.1^{+0.4}_{-0.4}$ & $1.9^{+4.2}_{-5.3}$ & $4.6^{+5.0}_{-3.3}$ \\
\bottomrule
\end{tabular}
\end{table}

Fig.~\ref{fig:stacking} shows the result. The averaged posterior means have $m \approx -0.4$ for all $N$, since averaging does not remove the shrinkage of the individual means (Appendix~\ref{app:shrinkage}). The stacked posterior, on the other hand, is already nearly unbiased for $N = 4$, and for all 32 galaxies gives $m_1 = -0.010 \pm 0.005$ and $m_2 = 0.001 \pm 0.004$, with additive biases consistent with zero at the level of a few $\times 10^{-4}$ (Table~\ref{tab:stacking}). This shows that the shear information in our NPE posteriors is nearly unbiased, even though their means are not, and that stacking them recovers it without any calibration.

These results assume a constant shear field, where every galaxy in a stacked cell is assumed to have the same shear. In real data, however, shear is not constant, so to validate, we simulated a cluster lensing scenario and stack galaxy likelihoods within spatial cells, which we will move to next.

\subsection{Likelihood Stacking in a Cluster Shear Field}
\label{ssec:cluster}

We now test likelihood stacking in a varying shear field. We simulate a dark matter halo of mass $M_\mathrm{200c} = 10^{15}\,M_\odot$, which follows an NFW profile~\citep{NFW97}, with a concentration from the mass--concentration relation of Duffy et al.~\citep{Duffy08}. We compute the shear $\boldsymbol{\gamma}$ and convergence $\kappa$ of the halo with \galsim~\citep{galsim}, and form the reduced shear $\mathbf{g} = \boldsymbol{\gamma}/(1-\kappa)$ at the position of each galaxy, assuming a Planck 2018 cosmology~\citep{Planck18}. Since our NPE is trained at a single redshift, we place all source galaxies at $z_s = 0.3$, and choose the lens redshift close to where the lensing signal of these sources peaks, at $z_l = 0.13$. The reduced shear of all our source galaxies stays below 0.075, within the prior range of our NPE.

We place the source galaxies uniformly within 3\,Mpc of the cluster center, drawing their number from a Poisson distribution with a mean density of 0.11\,arcmin$^{-2}$, which gives about 150 galaxies per cluster. This is the density of sample~3 of Xu et al.~\citep{KLDESI} over the DESI footprint, 0.14\,arcmin$^{-2}$ in our catalog of its targets, times the 75\% of these galaxies that lie behind the lens. Each galaxy takes the parameters of a randomly drawn Xu3 galaxy as in Sec.~\ref{ssec:test-set}, and we simulate its image and spectra and estimate its TF-weighted posterior as for the Xu3 test set. We repeat this for 50 independent realizations, each with new galaxies, positions, and noise, which will be used to bring down error bars in the bias analysis.

To stack the galaxies, we group them into cells and combine the likelihoods within each cell using Eq.~(\ref{eqn:stack}). Our cells consist of a central disk and four logarithmically spaced annuli out to 3\,Mpc, with each annulus divided such that all cells are roughly the same area and hold about 5 galaxies, which gives us 29 cells per cluster (Fig.~\ref{fig:cluster-cells}a). We rotate the likelihood of each galaxy into tangential and cross components $(g_\mathrm{t}, g_\times)$ about the cluster center before stacking.

Since the shear varies within a cell, and each shear estimate has a different posterior width, deriving a reference shear $\mathbf{g}_\mathrm{ref}$ with which to compare our measured shear $\hat{\mathbf{g}}$ with is non-trivial. For accuracy, we opt to use the full covariance of the shear likelihoods to derive the reference shear, rather than a simple galaxy weight. If the likelihood of each galaxy is approximately Gaussian with precision matrix $\mathbf{F}_i$ (the inverse of its covariance), the stacked shear converges to the reference shear
\begin{equation}
\label{eqn:gref}
    \mathbf{g}_\mathrm{ref} = \Big(\sum_{i} \mathbf{F}_i\Big)^{-1} \sum_{i} \mathbf{F}_i\, \mathbf{g}_i,
\end{equation}
the precision-weighted average of the true shears in the cell. Since KL constrains one direction of the shear much better than the other, $\mathbf{F}_i$ are strongly anisotropic, and $\mathbf{g}_\mathrm{ref}$ can differ noticeably from a plain average of the true shears or a scalar-weighted average.

Since the halo is axially symmetric, the true cross shear of every galaxy is zero by construction. A multiplicative bias needs a range of true shears to be measured against, so we measure both the multiplicative and the additive bias of the tangential shear, but only the additive bias of the cross shear. A bias in our likelihoods that leaks into the cross component would still show up as a non-zero stacked $g_\times$.

\begin{figure*}
    \includegraphics[width=\linewidth]{figs/cluster_cells.pdf}
    \caption{Likelihood stacking in the cells of the simulated cluster, for all 50 realizations. (a) Reference tangential shear $g_{\mathrm{ref,t}}$ of each cell (Eq.~\ref{eqn:gref}) and (b) stacked tangential shear $\hat{g}_\mathrm{t}$, both averaged over the realizations. The plus sign marks the cluster center. (c) Mean of $\hat{g} - g_\mathrm{ref}$ over all cells with $n$ galaxies, for $g_\mathrm{t}$ (blue circles) and $g_\times$ (red squares). Error bars are the standard errors of the mean. As for constant shear (Table~\ref{tab:stacking}), cells with only one or two galaxies are biased toward zero.}
    \label{fig:cluster-cells}
\end{figure*}

For each realization, we fit $\hat{g}_\mathrm{t} - g_{\mathrm{ref,t}} = m\, g_{\mathrm{ref,t}} + c_\mathrm{t}$ over its cells, weighting each cell by the inverse of its stacked posterior variance, and take the weighted mean of $\hat{g}_\times - g_{\mathrm{ref},\times}$ over its cells as $c_\times$. We then average each of these over the 50 realizations, and use the standard error of the mean as its uncertainty. We find $m = 0.06 \pm 0.05$, which did not increase significantly from our constant-shear test. Additive biases are $c_\mathrm{t} = (-4 \pm 6) \times 10^{-4}$, and $c_\times = (-2 \pm 3) \times 10^{-4}$, both consistent with zero. These uncertainties are larger than those for constant shear in Sec.~\ref{ssec:stacking}, since the tangential shear in a single cluster only spans about 0.01 to 0.04, and each cell holds only a few galaxies. \XW{Would be good to run more realizations here, to bring the uncertainties down to $\sim$1\% in $m$ and $\sim 10^{-4}$ in $c$, as for constant shear.} Likelihood stacking therefore shows no sign of additional bias in a non-constant shear field.

\subsection{Detecting Substructure}
\label{ssec:substructure}

Enabled by the reduced shape noise, one of the main advantages of KL is its ability to probe the substructure in galaxy clusters using individual galaxy shears, as~\citep{KLCluster} showed. Although the posterior means from individual galaxies from our NPE are biased (Appendix~\ref{app:posterior-means}), and the cells of Sec.~\ref{ssec:cluster} are too coarse to resolve a subhalo, we can demonstrate the ability to detect substructure with another method. Following the fitting method of Applegate et al.~\citep{WtG3}, we fit a model of the shear field directly to the likelihoods of all galaxies at once.

To test this, we add an NFW subhalo to the cluster of Sec.~\ref{ssec:cluster}. The subhalo has a mass of $M_\mathrm{200c} = 3\times10^{14}\,M_\odot$ (30\% of the main halo) and a concentration from Duffy et al.~\citep{Duffy08}, placed at the lens redshift and 1\,Mpc from the main halo center. We add the shear and convergence of the two halos before forming the reduced shear, which still stays within our training prior range for all galaxies. We do this for all 50 realizations of the main halo only field in Sec.~\ref{ssec:cluster}. The two versions (main-halo-only and main-with-subhalo) share the same galaxies, positions, and noise, so that any difference between them comes from the subhalo alone.

For a galaxy cluster model with parameters $\boldsymbol{\theta}$, which predicts the reduced shear $\mathbf{g}(\mathbf{x}_i; \boldsymbol{\theta})$ at the position $\mathbf{x}_i$ of each galaxy, the posterior of the parameters is
\begin{equation}
\label{eqn:halo-post}
    p(\boldsymbol{\theta} \mid \mathbf{c}_1, \dots, \mathbf{c}_N) \propto \pi(\boldsymbol{\theta}) \prod_{i=1}^{N} \mathcal{L}_i\big(\mathbf{g}(\mathbf{x}_i; \boldsymbol{\theta})\big),
\end{equation}
which is just Eq.~(\ref{eqn:stack}) with its argument replaced by the shear predicted using the sampled $\boldsymbol{\theta}$ at position $\mathbf{x}_i$. We fix the center of the main halo and take the concentrations of both halos from the mass--concentration relation, so that each halo is described fully by its mass and position. We fit three models: the main halo alone, the main halo and a subhalo at its true position, and the main halo and a subhalo whose position is left free. We use priors that are uniform in $\log_{10} M_\mathrm{200c}$, over $[13.5, 16]$ for the main halo and $[12, 15.5]$ for the subhalo, and a uniform prior on the subhalo position within $\pm 2.5$\,Mpc of the cluster center in each direction. We sample the posteriors of the models with a subhalo with \textsc{Nautilus}~\citep{Nautilus}, and evaluate the posterior of the main-halo-only model on a grid (since the only free parameter is the main halo mass).

To quantify the detection of the subhalo, we compare the maximum likelihood of the model with a subhalo at its true position, $\mathcal{L}_\mathrm{max,sub}$, with that of the main halo alone, $\mathcal{L}_\mathrm{max,main}$, through $\Delta\chi^2 = -2\ln(\mathcal{L}_\mathrm{max,main}/\mathcal{L}_\mathrm{max,sub})$, and quote the significance of the detection as $\sqrt{\Delta\chi^2}$. Fig.~\ref{fig:cluster-sigma} shows $\sqrt{\Delta\chi^2}$ for both versions of all realizations. With the subhalo, the median significance is $3.9\sigma$, above $2\sigma$ in 94\% and above $3\sigma$ in 76\% of the realizations. With the main halo only, the significance never exceeds $1.8\sigma$.

\begin{figure}
    \includegraphics[width=\linewidth]{figs/cluster_sigma.pdf}
    \caption{Detection significance $\sqrt{\Delta\chi^2}$ of the subhalo from the fits with the subhalo at its true position, for the 50 realizations with the subhalo (purple, solid) and with the main halo only (gray, dashed). The dotted line marks $2\sigma$.}
    \label{fig:cluster-sigma}
\end{figure}

The fits also recover the masses of the halos. On the data with the main halo only, the posterior medians of the main halo mass average to $\log_{10}(M_\mathrm{200c}/M_\odot) = 15.006 \pm 0.010$ over the realizations, consistent with the true value of 15. On the data with the subhalo, the posterior medians of the subhalo mass scatter around $\log_{10}(M_\mathrm{200c}/M_\odot) = 14.42$, close to the true value of 14.48, and the 68\% credible interval contains the true mass in 80\% of the realizations.

When the position of the subhalo is also left free, the fit still locates it. Fig.~\ref{fig:cluster-localization} shows the posterior median of the subhalo position in each realization. It lies within 0.5\,Mpc of the true position in 82\% of the realizations, with a median offset of 0.22\,Mpc.

\begin{figure}
    \includegraphics[width=0.8\linewidth]{figs/cluster_localization.pdf}
    \caption{Posterior median of the subhalo position from the fits with the subhalo position left free, for the 50 realizations with the subhalo (purple points). The plus sign marks the center of the main halo, the cross the true position of the subhalo, and the circle a radius of 0.5\,Mpc around it.}
    \label{fig:cluster-localization}
\end{figure}

Overall, by fitting a model of the shear field to the NPE likelihoods, a subhalo with 30\% of the mass of the main halo can be detected in a single cluster at the source density of a DESI-like KL sample, and located to within about 0.5\,Mpc, without any calibration of the NPE.
